% ============================================================
%  第 2 章  相关技术与理论基础
% ============================================================

\chapter{相关技术与理论基础}

% ================================================================
\section{Apache Kafka 消息队列}

\subsection{Kafka 基本架构}

Apache Kafka~\cite{kafka2011} 是由 LinkedIn 公司开发、后捐献给 Apache 基金会的分布式流式消息系统，基于发布/订阅（Publish/Subscribe）模式设计。Kafka 的核心概念包括：

\begin{itemize}
  \item \textbf{Broker}：Kafka 集群的服务节点，负责消息的持久化存储与转发。
  \item \textbf{Topic}：消息的逻辑分类，每个 Topic 可配置多个 Partition 实现并行读写。
  \item \textbf{Producer}：向指定 Topic 写入消息的客户端进程。
  \item \textbf{Consumer Group}：消费者组内各实例分配不同 Partition，保证每条消息在组内仅被消费一次，实现水平扩展。
  \item \textbf{Offset}：Consumer 在 Partition 中消费进度的游标，由 Consumer 自主管理，支持精确一次（Exactly-Once）或至少一次（At-Least-Once）语义。
\end{itemize}

\subsection{Kafka 的核心特性}

Kafka 在分布式流式数据传输领域具备以下关键特性：

\begin{enumerate}
  \item \textbf{高吞吐、低延迟：}Kafka 采用顺序磁盘 I/O 与零拷贝（Zero-Copy）技术，在普通硬件上可实现百万级消息/秒的吞吐量，99\% 分位延迟通常低于 10~ms。
  \item \textbf{持久化存储：}消息以 Log 文件形式持久化到磁盘，默认保留时间可配置（本系统设置为 168 小时），支持历史消息回溯消费。
  \item \textbf{水平扩展：}通过增加 Partition 数和 Consumer 实例，吞吐量可线性扩展。
  \item \textbf{精确一次语义：}结合 Idempotent Producer 与 Transactional API，支持跨 Partition 原子写入，可配合 Spark Structured Streaming 实现端到端 Exactly-Once 语义。
\end{enumerate}

\subsection{本系统中 Kafka 的应用}

本系统使用 Kafka 作为爬虫与 Spark 流处理管道之间的解耦缓冲层。爬虫将采集到的评论序列化为 JSON 格式后写入 Topic \code{public\_opinion\_raw}，Spark Structured Streaming 以 Kafka Source 方式订阅该 Topic，按设定的微批次周期（60~秒）消费消息。这一设计将数据生产（爬虫）与数据消费（Spark）完全解耦，任一端发生临时故障均不影响对方，Kafka 的消息持久化也保障了 Spark 异常重启后可从上次提交的 Offset 继续消费，不产生数据丢失。

% ================================================================
\section{Apache Spark Structured Streaming}

\subsection{Structured Streaming 的设计哲学}

Apache Spark Structured Streaming~\cite{armbrust2018structured} 是 Spark 2.0 引入的流处理引擎，其核心设计理念是将\textbf{无界流数据}抽象为一张持续追加行的无界关系表（Unbounded Table）。用户在此表上编写与静态 DataFrame 完全一致的 SQL/DSL 查询，引擎负责将其转化为增量的流式执行计划。这一设计使得流式程序与批处理程序可共用同一套 API，大幅降低了开发复杂度。

\subsection{微批次处理模型}

Structured Streaming 支持三种触发模式：固定周期（Fixed Interval）微批次、一次性触发（Once）和连续处理（Continuous Processing）。本系统采用固定周期微批次，触发间隔为 60~秒，即系统每 60~秒从 Kafka 拉取一次消息作为 DataFrame，通过 \code{foreachBatch} 接口传递给自定义处理函数。微批次模型的优势在于：（1）端到端 Exactly-Once 语义实现简单；（2）可借助 Spark 批处理优化器对每批次 SQL 计划进行完整优化；（3）与 Delta Lake ACID 事务写入天然兼容。

\subsection{\texttt{foreachBatch} 接口}

\code{foreachBatch} 是 Structured Streaming 提供的核心扩展点，允许用户为每个微批次提供一个 \code{(DataFrame, batchId) -> Unit} 形式的回调函数。该接口突破了原生 Sink 的局限性，使得批次内部可以执行任意复杂的逻辑，包括：调用外部 Python 库（如本系统中的 StructBERT 推理）、执行多表 Merge 操作、发送 HTTP 告警请求等。本系统的核心业务逻辑（情感推理、热度计算、Delta Lake 三层写入、告警检测）均在 \code{process\_batch} 回调函数内实现，如第~4~章所述。

\subsection{Checkpoint 与容错保障}

Structured Streaming 通过 Checkpoint 机制实现流处理的容错性。Checkpoint 目录存储了三类关键信息：（1）Kafka Offset 提交记录，记录每批次消费的消息范围；（2）执行计划分析状态；（3）聚合算子的状态存储（State Store）。系统异常重启时，引擎从 Checkpoint 恢复上次提交的 Offset，保证从断点处续算，避免重复处理或遗漏消息，从而为下游 Delta Lake 写入提供 At-Least-Once 保障，结合 Silver 层的 Merge 去重操作，整体实现 Exactly-Once 语义。

% ================================================================
\section{Delta Lake 数据湖仓架构}

\subsection{Delta Lake 概述}

Delta Lake~\cite{armbrust2020delta} 是开源的、基于 Parquet 文件格式的存储层，通过在数据目录中维护一个 JSON 格式的事务日志（Transaction Log，即 \code{\_delta\_log} 目录）在对象存储之上实现了 ACID 事务语义。每次写操作（\code{INSERT}、\code{UPDATE}、\code{DELETE}、\code{MERGE}）均原子地追加一条日志记录；读操作通过重放日志构建表的最新一致性快照，从而解决了 Parquet 文件本身不支持事务的根本性缺陷。

\subsection{核心特性}

\begin{itemize}
  \item \textbf{ACID 事务：}通过乐观并发控制（OCC）实现 Serializable 级别隔离，保障并发读写的一致性。
  \item \textbf{Schema 强制与演化：}写入时自动校验数据 Schema，防止脏数据；支持在不中断读取的情况下安全地新增字段（Schema Evolution）。
  \item \textbf{Time Travel：}通过事务日志支持按版本号或时间戳回溯任意历史快照，便于数据审计与回滚。
  \item \textbf{流批统一：}Structured Streaming 可以直接使用 \code{writeStream.format("delta")} 流式写入，也可用 \code{spark.read.format("delta")} 批量读取，同一张表支持流读和批读并存。
  \item \textbf{MERGE INTO：}支持 SQL \code{MERGE INTO} 语句，可在单次操作中同时处理插入、更新和删除逻辑，是本系统 Silver 层去重精炼的核心操作。
\end{itemize}

\subsection{三层数据架构设计思想（Bronze-Silver-Gold）}

Medallion 架构~\cite{databricks2021medallion}（又称三层数据湖架构）是 Databricks 推荐的 Delta Lake 数据组织模式，将数据按加工程度划分为三个层次：

\begin{itemize}
  \item \textbf{Bronze 层（原始层）：}原样保存从消息队列接收的原始数据，仅做最少的格式转换（如字符串解析）。数据以 Append 模式写入，保留完整历史，用于数据回溯与问题排查。
  \item \textbf{Silver 层（精炼层）：}对 Bronze 层数据进行去重（基于业务主键 \code{event\_id}）、字段校验、特征富化（情感标签、热度分数、风险等级）等处理后，以 Merge 模式写入，形成面向分析的宽表。
  \item \textbf{Gold 层（应用层）：}基于 Silver 层数据进行多粒度（按小时、按天）聚合，生成直接面向前端 API 消费的统计指标表，以全量 Overwrite 模式保持幂等性。
\end{itemize}

三层架构的核心价值在于将"原始数据保存"与"数据加工结果"分层隔离，上层写入失败不影响底层数据完整性，同时为不同消费场景（实时查询 vs 历史回溯）提供了针对性优化的数据视图。

% ================================================================
\section{StructBERT 中文情感分析模型}

\subsection{BERT 预训练语言模型基础}

BERT（Bidirectional Encoder Representations from Transformers）~\cite{devlin2019bert} 由 Google 于 2018 年提出，采用 Transformer Encoder 架构，通过掩码语言模型（Masked Language Modeling, MLM）与下一句预测（Next Sentence Prediction, NSP）两个无监督预训练任务，在海量通用语料上学习深层双向语言表示。预训练完成后，通过在顶层添加任务相关分类头并进行少量有监督微调（Fine-tuning），BERT 在 GLUE 基准的情感分类、文本蕴含等 11 项下游任务上全面超越了当时的 SOTA，引发了自然语言处理领域的"预训练大模型"时代。

\subsection{StructBERT 的结构增强}

StructBERT~\cite{wang2020structbert} 是阿里达摩院在 BERT 基础上的改进版本，针对 BERT 原始预训练任务设计中的不足提出了两项结构感知增强：

\begin{enumerate}
  \item \textbf{词序重建任务（Word Structural Objective）：}在 MLM 基础上，额外对部分词片段进行随机打乱，要求模型在预测被掩盖词的同时还原词序，迫使模型学习更强的词级局部句法依存关系。
  \item \textbf{句子顺序预测任务（Sentence Structural Objective）：}将原始 NSP 任务（判断两句话是否相邻）替换为三分类任务（判断两句话的顺序关系：正序、逆序或随机），提供了更细粒度的篇章级语义监督。
\end{enumerate}

上述两类增强任务使 StructBERT 在 BERT-Large 同等参数量下，在 STS-B（语义相似度）、MNLI（自然语言推断）等多项 GLUE 子任务上取得了显著提升~\cite{wang2020structbert}。阿里达摩院基于 StructBERT 进一步在中文情感分类语料上进行了任务特定微调，发布了 ModelScope 平台上的 \code{iic/nlp\_structbert\_sentiment-classification\_chinese-base} 模型，该模型即本系统所采用的情感分析推理引擎。

\subsection{本系统中的情感分析设计}

本系统通过 ModelScope Pipeline 接口调用 StructBERT，输入单条或批量中文文本，输出情感标签（\code{positive} / \code{negative} / \code{neutral}）及对应置信度分数。具体工程设计如下：

\begin{enumerate}
  \item \textbf{懒初始化单例：}在 \code{foreachBatch} 首次触发时完成模型权重加载（约 400~MB），后续批次复用同一实例，避免每批次重复加载带来的数秒延迟。
  \item \textbf{批量推理：}将批次内所有评论文本分块（默认每块 32 条）后批量送入模型，充分利用并行计算，相较逐条推理吞吐量提升约 5 倍。
  \item \textbf{中性判定阈值：}当模型对 \code{positive} 与 \code{negative} 两类别的置信度均低于可配置阈值（默认 0.72）时，将情感标签归为 \code{neutral}，减少边界样本的强制归类误差。
  \item \textbf{异常降级：}推理过程中若发生任何未捕获异常，将该条评论的情感标签默认为 \code{neutral}，热度加成取中性值 0.5，保障管道不中断。
\end{enumerate}

% ================================================================
\section{FastAPI 后端框架}

FastAPI~\cite{fastapi2019} 是基于 Python 3.6+ 类型注解的现代高性能 Web 框架，底层使用 ASGI 服务器（Uvicorn）实现异步请求处理。FastAPI 的核心优势包括：（1）基于 Pydantic 的请求/响应模型自动验证与序列化；（2）自动生成符合 OpenAPI 3.0 规范的交互式文档（Swagger UI）；（3）原生支持 \code{async/await}，与异步 I/O 生态无缝集成；（4）与 Python 数据生态（Pandas、NumPy）深度兼容，可直接在接口函数中进行 DataFrame 操作后返回 JSON。

在本系统中，FastAPI 负责将 Delta Lake 表数据封装为 RESTful 接口，供前端 Vue 3 大屏消费。由于 \code{deltalake} Python 库（\code{delta-rs}）支持无需 JVM 直接读取 Delta 表的 Parquet 文件，因此 FastAPI 服务可以完全独立于 Spark 运行，显著降低了部署复杂度与资源开销。

% ================================================================
\section{Vue 3 与 ECharts 可视化技术}

\subsection{Vue 3 前端框架}

Vue 3~\cite{vue3} 是尤雨溪主导的渐进式前端框架，2020 年发布的 3.0 版本带来了以下重要改进：（1）基于 Proxy 的响应式系统，替代 Vue 2 的 \code{Object.defineProperty}，支持对数组和新增属性的自动追踪；（2）Composition API，以 \code{setup()} 函数替代 Options API，实现更细粒度的逻辑复用与代码组织；（3）更优的 TypeScript 支持；（4）更小的运行时体积（约 10~KB gzip）。

本系统前端基于 Vue 3 Composition API 构建单页应用（SPA），利用 \code{ref}、\code{computed} 与 \code{watch} 管理应用状态，通过原生 \code{fetch} API 与后端通信，整体代码组织于单一 \code{App.vue} 文件中，保持了结构的简洁性。

\subsection{Apache ECharts 数据可视化}

Apache ECharts~\cite{echarts} 是百度开源、后捐献给 Apache 基金会的高性能数据可视化库，基于 Canvas 与 SVG 双渲染引擎，支持折线图、柱状图、饼图、热力图、地图等数十种图表类型，具备自动适配移动端的响应式能力。ECharts 的 \code{setOption} 接口采用声明式配置模式，极大简化了图表的动态更新逻辑：只需在数据刷新时重新调用 \code{setOption}，ECharts 会自动完成数据绑定与动画过渡。本系统使用 ECharts 绘制评论数/负面数趋势折线图、负面占比趋势折线图与热度趋势折线图，三图联动响应话题切换与粒度切换操作。

\subsection{技术栈选型小结}

本系统在多个技术决策点进行了方案对比与选型，关键选型决策总结如表~\ref{tab:tech-selection}~所示。

\begin{table}[htbp]
\centering
\caption{关键技术选型对比}
\label{tab:tech-selection}
\begin{tabularx}{\textwidth}{@{}p{2cm}p{3.6cm}p{4cm}X@{}}
\toprule
\tablehead{技术模块} & \tablehead{候选方案 A（未选）} & \tablehead{候选方案 B（已选）} & \tablehead{选型理由} \\
\midrule
消息队列    & RabbitMQ               & Apache Kafka                  & Kafka 持久化、回溯消费、Spark 原生集成 \\
流处理引擎  & Apache Flink           & Spark Structured Streaming    & 与 Python/PyTorch 生态集成更成熟 \\
存储层      & PostgreSQL / MySQL     & Delta Lake (Parquet)           & 流批统一、ACID 湖仓、Schema 演化 \\
情感模型    & 词典规则（HowNet）     & StructBERT                    & 语义理解能力强、ModelScope 开箱即用 \\
后端框架    & Flask                  & FastAPI                       & 异步支持、自动校验、OpenAPI 文档 \\
前端框架    & React                  & Vue 3                         & 轻量、上手快、Composition API 灵活 \\
数据可视化  & D3.js                  & Apache ECharts                & 声明式配置、移动端适配、中文社区丰富 \\
\bottomrule
\end{tabularx}
\end{table}
